% III. Проектиране на системата.
% Слоеве, отговорности, граници между тях, формат на индексите, поведение при
% грешка. Числа за производителност тук няма: те се появяват едва в Раздел VI.
\section{Проектиране на системата}
\label{sec:design}

\subsection{Слоеве и отговорности}

Системата е разделена на слоеве, като всеки от тях зависи само от слоевете под себе
си и общува с тях през тесен договор. Слоят на синтаксиса чете Turtle и N-Triples и
превръща текста в поток от триплети над термини. Речникът на термините съпоставя
всеки термин с числов идентификатор и обратно. Слоят на индексите поддържа трите
пермутирани подредби на кодираните триплети. Анализаторът на SPARQL превръща текста
на заявката в алгебрично дърво. Планировчикът преобразува това дърво в дърво от
оператори, като избира реда на съединенията и индекса за всеки шаблон. Механизмът за
изпълнение обхожда операторите и произвежда решенията. Разположението на слоевете и
посоката на данните са показани на \figref{fig:architecture}.

\begin{figure}[H]
\centering
\begin{tikzpicture}[node distance=0.8cm and 1.6cm]
  \node[block] (syntax) {Синтаксис\\Turtle, N-Triples};
  \node[block, below=of syntax] (dict) {Речник на термините};
  \node[block, below=of dict] (idx) {Пермутирани индекси\\SPO, POS, OSP};
  \node[block, right=of syntax] (sparql) {Анализатор на SPARQL\\алгебрично дърво};
  \node[block, below=of sparql] (plan) {Планировчик\\ред на съединенията};
  \node[block, below=of plan] (exec) {Механизъм за изпълнение};
  \draw[line] (syntax) -- (dict);
  \draw[line] (dict) -- (idx);
  \draw[line] (sparql) -- (plan);
  \draw[line] (plan) -- (exec);
  \draw[line] (exec) -- (idx);
\end{tikzpicture}
\caption{Слоеве на системата. Лявата колона е пътят на данните при зареждане, дясната е пътят на заявката. Механизмът за изпълнение достига данните само през индексите, а планировчикът чете от тях единствено статистики}
\label{fig:architecture}
\end{figure}

\subsection{Речник на термините}

Речникът поддържа две посоки. Посоката от термин към идентификатор се използва при
зареждане и при превръщане на константите в заявката, а обратната само при извеждане
на резултата. Ключът на съпоставянето е записът на термина в N-Triples, тъй като този
запис е каноничен: един и същ термин, написан в Turtle с префикс или изцяло, дава
един и същ ключ и се вписва в речника само веднъж.

Идентификаторът е 64 бита, разделени на две части. Старшите 8 бита носят вида на
термина, тоест IRI, празен възел или литерал. Младшите 56 бита носят поредния номер,
раздаван от речника. Това разделение позволява на механизма за изпълнение да отговори
на \code{isIRI}, \code{isLiteral} и \code{isBlank} без да чете самия термин, и оставя
място за над $7 \cdot 10^{16}$ различни термина, което е далеч над обхвата на вграден
случай на употреба. Стойността нула е запазена и означава „няма свързване“, затова
поредните номера започват от единица.

Празните възли са локални за документа, от който идват. При зареждане етикетът им се
предхожда от име на обхват, извлечено от името на файла, така че два документа, заредени
в едно хранилище, не могат да се сблъскат. Анонимните празни възли, тоест написаните
като \code{[]} или получени от съкратен списък, получават генерирано име в същия обхват.

\subsection{Кодирани триплети и пермутирани индекси}

След кодиране триплетът е три числа с фиксиран размер, тоест 24 байта. Всеки индекс
е редица от такива триплети, подредена лексикографски по съответната пермутация.
Поддържат се три пермутации: подлог, сказуемо, допълнение (SPO); сказуемо,
допълнение, подлог (POS); и допълнение, подлог, сказуемо (OSP).

Редиците се държат изцяло в оперативната памет. Отразяването им във файлове и
компресията на индексите остават извън обхвата по причината, дадена в
Раздел~\ref{sec:alternatives}, и са записани като посока за по-нататъшна работа.

Търсенето по шаблон се свежда до двоично търсене на границите на префикса,
съответстващ на свързаните позиции, последвано от последователно четене на диапазона.
Тази форма е избрана заради две свойства: последователното четене е предвидимо за
подсистемата на паметта, а подредеността позволява съединяването на два входа да бъде
сливане на сортирани редици, без изграждане на хеш таблица.

\subsection{Покритие на шаблоните от трите индекса}

Твърдението, върху което стои изборът на точно три пермутации, е следното: за всяка
от осемте комбинации от свързани и свободни позиции съществува индекс, в който всички
свързани позиции образуват префикс. \tabref{tab:index-choice} изброява осемте случая
заедно с избрания индекс и дължината на префикса. Проверката на това твърдение е
записана като тест, а не оставена като разсъждение.

\begin{table}[H]
\caption{Избор на индекс според свързаните позиции на шаблона}
\label{tab:index-choice}
\centering
\fontsize{10}{12}\selectfont
\begin{tabular}{@{}L{3.4cm}L{1.6cm}L{2.0cm}L{3.9cm}@{}}
\toprule
\textbf{Свързани позиции} & \textbf{Индекс} & \textbf{Префикс} & \textbf{Подредба на изхода} \\
\midrule
няма             & SPO & 0 & по подлог \\
подлог           & SPO & 1 & по сказуемо \\
сказуемо         & POS & 1 & по допълнение \\
допълнение       & OSP & 1 & по подлог \\
подлог, сказуемо & SPO & 2 & по допълнение \\
сказуемо, допълнение & POS & 2 & по подлог \\
подлог, допълнение & OSP & 2 & по сказуемо \\
и трите          & SPO & 3 & единичен резултат \\
\bottomrule
\end{tabular}
\end{table}

Последната колона е също толкова съществена, колкото и втората. Обхождането на
диапазон в даден индекс произвежда триплетите подредени по следващата позиция от
пермутацията, тоест по позицията непосредствено след префикса. Тази подредба е
безплатна и е единственото условие, при което съединяването чрез сливане е допустимо.

\subsection{Алгебрично дърво и план за изпълнение}

Анализаторът на SPARQL произвежда дърво, чиито възли съответстват на операторите от
алгебрата на спецификацията: базов графов шаблон, съединение, незадължително
съединение, обединение, филтър, групиране с агрегация, проекция, премахване на
повторенията, подреждане и ограничение на броя решения~\cite{w3c_sparql11_query}.
Планировчикът работи само върху това дърво и не познава синтаксиса, което позволява
двете части да се тестват поотделно.

Планът за изпълнение е второ дърво, този път от оператори. Преобразуването прави три
неща. Първо, изброява променливите на цялата заявка и дава на всяка номер на позиция,
така че по време на изпълнение свързването на променлива е достъп по индекс в масив, а
не търсене по име. Второ, кодира константите на всеки шаблон през речника; константа,
която речникът не познава, прави шаблона празен и това се вижда преди каквото и да е
обхождане. Трето, фиксира реда на съединенията и избира индекс и стратегия за всяко
съединение.

\subsection{Оценка на мощността и ред на съединенията}

Под \term{мощност} на шаблон се разбира броят на триплетите, които го напасват при
свързани само неговите константи. В избраната схема на индексите тази величина не се
оценява приблизително, а се пресмята точно: тя е дължината на диапазона, който двоичното
търсене връща, и струва $O(\log n)$. Планировчикът работи от точни числа, което
премахва цял клас грешки, произтичащи от разминаване между оценка и действителност.

Редът на съединенията се избира лакомо. Първи се взема шаблонът с най-малка мощност.
На всяка следваща стъпка се взема шаблонът с най-малка мощност измежду онези, които
споделят променлива с вече избраните; ако такъв няма, се взема най-малкият измежду
останалите. Предпочитанието към свързаност е важно, защото шаблон без обща променлива
дава декартово произведение, чийто размер расте като произведение на мощностите.

За да може ползата от планирането да бъде измерена, а не предположена, планировчикът
приема флаг, който запазва реда на изписване на заявката. Двата режима произвеждат едни
и същи решения и се различават само по време на изпълнение.

\subsection{Оператори за изпълнение}

Изпълнението е верига от итератори. Всеки оператор поддържа две операции: отваряне,
което получава вече направените свързвания от обхващащия оператор, и извличане на
следващото решение. Този договор е причината съединяването чрез вложени цикли да бъде
съединяване по индекс: вътрешният оператор получава свързванията на външния и стеснява
диапазона си с тях, вместо да чете всичко и да отсява.

\begin{table}[H]
\caption{Оператори на механизма за изпълнение}
\label{tab:operators}
\centering
\fontsize{10}{12}\selectfont
\begin{tabular}{@{}L{4.4cm}L{3.3cm}L{5.1cm}@{}}
\toprule
\textbf{Оператор} & \textbf{Материализира} & \textbf{Забележка} \\
\midrule
Обхождане по индекс & не & избира индекс по свързаните позиции \\
Съединяване чрез вложени цикли по индекс & не & отваря дясната страна за всеки ред отляво \\
Съединяване чрез сливане & частично & буферира само групата с равен ключ \\
Незадължително съединяване & не & при липса на съвпадение връща реда отляво \\
Обединение & не & изчерпва левия вход, после десния \\
Филтър & не & изчислява израза върху всеки ред \\
Проекция & не & изчиства позициите извън проекцията \\
Премахване на повторения & да, ключовете & пази множество от вече върнати редове \\
Подреждане & да & сортиращите ключове се пресмятат веднъж на ред \\
Групиране с агрегация & да & никоя група не е завършена преди края на входа \\
Ограничение и отместване & не & спира входа рано \\
\bottomrule
\end{tabular}
\end{table}

Материализиране се извършва само там, където то е присъщо на оператора. Подреждането
не може да върне първия ред, преди да е видяло последния, а агрегацията не може да
затвори група по същата причина. Останалите оператори работят поточно, което е и
причината ограничението на броя решения да спира работата рано.

Съединяването чрез сливане е допустимо само когато и двата входа са обхождания по
индекс, подредени по общата променлива. Според \tabref{tab:index-choice} това се случва
при шаблони със свързано сказуемо, които споделят допълнението си: и двата се четат от
POS и излизат подредени по допълнение. В план с лява дълбочина условието може да бъде
изпълнено само за първата двойка, тъй като изходът на съединение вече не е подреден.

\subsection{Слой за извеждане по RDFS}

Извеждането по RDFS е реализирано като материализация: правилата се прилагат след
зареждане и получените триплети се вмъкват в индексите. Реализираното подмножество е
следното, с номерата на правилата от спецификацията~\cite{w3c_rdf_schema}:

\begin{itemize}[topsep=0pt,itemsep=0pt,leftmargin=*]
\item \textbf{rdfs2}: ако свойство има област $C$ и съществува триплет с това свойство,
      подлогът му е от тип $C$;
\item \textbf{rdfs3}: ако свойство има обхват $C$, допълнението е от тип $C$;
\item \textbf{rdfs5}: подсвойството е транзитивно;
\item \textbf{rdfs7}: триплет по подсвойство важи и по надсвойството;
\item \textbf{rdfs9}: ресурс от подклас е и от надкласа;
\item \textbf{rdfs11}: подкласът е транзитивен.
\end{itemize}

Правилата се прилагат до неподвижна точка. Итерирането е необходимо, защото правилата
се подхранват едно друго: обхватът дава тип, а подкласът след това разширява този тип.
Броят на обиколките е ограничен отгоре, за да не може дефект в правилата да доведе до
безкраен цикъл вместо до съобщение за грешка.

Извеждането по време на заявка не е реализирано. То, аксиоматичните триплети на RDFS и
правилата за членовете на контейнери са съзнателно извън обхвата.

\subsection{Поведение при грешка}

Грешките се разделят на три вида: синтактична грешка във входния документ или в
заявката, нарушено съответствие в самото хранилище и изчерпан ресурс. Първият вид се
съобщава с позиция в текста, ред и колона. Зареждането е двустъпково: триплетите се
натрупват в междинен буфер и се прехвърлят в хранилището едва след като целият документ
е прочетен, така че документ с грешка по средата не оставя частично заредени данни.
Вторият вид спира операцията, защото продължаването върху повреден индекс дава тихо
грешен отговор, което е по-лошо от отказ. Обхождане на хранилище, чиито индекси не са
изградени след последната промяна, попада именно тук и се отхвърля.

Конструкция на SPARQL извън поддържаното подмножество се отхвърля с изрично съобщение.
Пренебрегната конструкция би дала заявка, която означава нещо различно от написаното,
без никакъв признак за това.
